\documentclass[oneside,a4paper,12pt]{article} % Gibt an: Papierformat, Schriftgröße

\usepackage{thesis}

% Hier werden die Abkürzungen definiert. Sofern ein Abkürzungsverzeichnis verwendet wird bitte entkommentieren.
%\newacronym{aifb}{AIFB}{Institut für Angewandte Informatik und Formale Beschreibungsverfahren}
\newacronym{kit}{KIT}{Karlsruher Institut für Technologie}
\newacronym[plural={Knowledge Graphs}]{kg}{KG}{knowledge graph}
\newacronym{bpe}{BPE}{byte pair encoding}
\newacronym{test}{TEST}{test}
\newacronym[plural={large language models}]{llm}{LLM}{large language model}
\newacronym{nlp}{NLP}{natural language processing}
\newacronym{kgqa}{KGQA}{knowledge graph question answering}
\newacronym[plural={natural language questions}]{nlq}{NLQ}{natural language question}
\newacronym[plural={pre-trained language models}]{plm}{PLM}{pre-trained language model}
\newacronym{lm}{LM}{language modeling}
\newacronym{nlu}{NLU}{natural language understanding}
\newacronym{nlg}{NLG}{natural language generation}


\begin{document}
\setlanguageGerman % Sprache einstellen.

% Hier kommt der ganze Vorspann. Bei Verwendung von Abkürzungs-, Abbildungs- oder Tabellenverzeichnisse bitte in dieser Datei entsprechend entkommentieren.
\pagenumbering{roman}
\begin{titlepage}
% font / Schriftart
%------------------	


\begin{figure}[htbp]
  %\centering
  \begin{minipage}[b]{7cm}
    \ifthenelse{\boolean{english}}{\includegraphics[height=1.8cm]{images/kit_logo_en}}{\includegraphics[height=1.8cm]{images/kit_logo_de}}
  \end{minipage}
  \begin{minipage}[b]{8.5cm}
	\begin{flushright}


    \includegraphics[height=1.8cm]{images/aifb_logo}  
		\end{flushright}
  \end{minipage}
\end{figure}

\vspace*{2.5cm}
\begin{center}
		%TODO Titel angeben
		{\textbf{\Huge \ifthenelse{\boolean{english}}{Title of the Thesis}{Titel der Arbeit}}}
		\vspace*{1cm}\\
		%TODO Art der Arbeit
		{\Large \ifthenelse{\boolean{english}}{Seminar Paper / Bachelor Thesis / Master Thesis}{Seminararbeit / Bachelorarbeit / Masterarbeit}}\\[1cm]
		{\Large \ifthenelse{\boolean{english}}{by}{von}}\\[1cm]
		%\vspace*{1cm}
		%TODO Name des Studenten
		{\huge Vorname Nachname}\\[0.5cm]
		%TODO Studiengang des Studenten
		{\Large \ifthenelse{\boolean{english}}{Degree Course: Industrial Engineering and Management B.Sc.}{Studiengang: Wirtschaftsingenieurwesen}}\\[0.25cm]
		%TODO MatNr angeben
		{\Large \ifthenelse{\boolean{english}}{Matriculation Number}{Matrikelnummer}: 1234567}\\
		\vspace*{1.5cm}
		{\Large\ifthenelse{\boolean{english}}{
			Institute of Applied Informatics and Formal Description Methods (AIFB)
			\\[0.5cm]
			KIT Department of Economics and Management}{
			Institut für Angewandte Informatik und Formale Beschreibungsverfahren (AIFB)
			\\[0.5cm]
			KIT-Fakultät für Wirtschaftswissenschaften}
		}
	\end{center}
	\vspace*{1.5cm}
\Large{
\begin{center}
\begin{tabular}[ht]{l c l}
	
  %TODO Referent/-in angebenn
  %Referent/-in sind die Professoren, die die Arbeit bewerten.
  \ifthenelse{\boolean{english}}{Advisor}{Prüfer}: & \hfill  & *Titel Vor- und Nachname* \\
  %Referent/-in sind die Professoren, die die Arbeit bewerten.
  \ifthenelse{\boolean{english}}{Second Advisor}{Zweiter Prüfer}: & \hfill  & *Titel Vor- und Nachname* \\
  %TODO Betreuer/-in angeben
  \ifthenelse{\boolean{english}}{Supervisor}{Betreuer}: & \hfill  & *Titel Vor- und Nachname* \\
  %TODO Abgabedatum angeben
  \ifthenelse{\boolean{english}}{Submitted}{Eingereicht am}: & \hfill  & \today\\
 
\end{tabular}
\end{center}

}


%\include{titel}
  \vfill
	
	\small{ KIT -- \ifthenelse{\boolean{english}}{The Research University in the Helmholtz Association}{Die Forschungsuniversität in der Helmholtz-Gemeinschaft}
} \hfill
	\small{\textbf{\url{www.kit.edu}} }
\end{titlepage}
%Für den Fall eines Vorworts entkommentieren:
%\mbox{}\thispagestyle{empty}

\vspace*{1cm}

{\Large \textbf{\ifthenelse{\boolean{english}}{Prologue}{Vorwort}}} 

\bigskip

\ifthenelse{\boolean{english}}
{\selectlanguage{english}
In English theses, German prologue is probably to be used, so paste it here. \blindtext
}
{\selectlanguage{ngerman}
Hier könnte ein Vorwort stehen.
}
\thispagestyle{empty}\section*{\ifthenelse{\boolean{english}}{Abstract}{Zusammenfassung}}
An dieser Stelle erfolgt eine knappe Zusammenfassung der vorliegenden Arbeit ([engl.] Abstract), die maximal ca. 200 Worte umfassen sollte. Der Sinn und Zweck dieser Zusammenfassung liegt darin, einem interessierten Leser die Entscheidung zu erleichtern, die vorliegende Arbeit überhaupt zu lesen bzw. vor dem Lesen der Arbeit erst einmal in Erfahrung zu bringen, worum es dabei geht. Also eine knappe, motivierende Hinführung zum Problem und wie Sie es gelöst haben.

Wenn Sie eine Zusammenfassung schreiben, bedenken Sie, dass diese oft auch alleine publiziert wird, d.h. sie sollte unabhängig vom nachfolgenden explizit dargestellten Inhalt der Arbeit für den Leser verständlich sein. Daher ist es immer sinnvoll, diese Zusammenfassung erst ganz am Ende zu schreiben, wenn Sie die eigentliche Arbeit bereits abgeschlossen haben.
%\setcounter{page}{3}
\tableofcontents\newcounter{roemisch} %\setcounter{roemisch}{\value{page}}
% Für den Fall von Abkürzungs-, Abbildungs- oder Tabellenverzeichnissen entkommentieren
%\ifthenelse{\boolean{english}}
%{\cleardoublepage\printnoidxglossary[type=acronym,sort=letter,nonumberlist,title={List of Abbreviations}]}
%{\cleardoublepage\printnoidxglossary[type=acronym,sort=letter,nonumberlist,title={Abkürzungsverzeichnis}]}
%\cleardoublepage\listoffigures
%\cleardoublepage\listoftables
%\cleardoublepage\setcounter{page}{2} 
\cleardoublepage
\pagenumbering{arabic}

% Füge hier die Kapitel per Verweis ein. Der Befehl \include fügt die angegebene tex-Datei an der jeweiligen Stelle ein. Die eingefügte Datei wird als normaler Teil des Quelltextes mitverarbeitet und ist daher LaTeX-Code (ohne Vorspann und \begin{document}...\end{document}).
\section{Überschrift1}
\label{Label1}

Hier steht Text. Diese Arbeit zitiert eine Beispielpublikation \cite{example_reference}.

% Beispiel für die Verwendung eines Abkürzungsverzeichnisses. Bitte entkommentieren:
%Das \gls{aifb} gehört zum \gls{kit}.

\subsection{Überschrift2}
\label{Label2}

Hier steht Text.

\section{Überschrift1}
\label{Label1}

Hier steht Text.

\subsection{Überschrift2}
\label{Label2}

Hier steht Text.

\section{Überschrift1}
\label{Label1}

Hier steht Text.

\subsection{Überschrift2}
\label{Label2}

Hier steht Text.

% ...usw.



% Anhang / Appendix
\appendix % Ab hier wird mit A, B, ... weiternummeriert.
% Für den Fall eines Anhangs entkommentieren
%\section{Some Appendix}
\subsection{Some subsection}


% Literaturverzeichnis
\bibliographystyle{acm}
\bibliography{YOUR_thesis} % Datei mit Literaturangaben einbinden

% Schriftliche Erklärung
\newpage
\mbox{}\thispagestyle{empty}

\vspace*{1cm}

{\Large \textbf{\ifthenelse{\boolean{english}}{Assertion}{Erklärung}}} 

\bigskip

\ifthenelse{\boolean{english}}
{\selectlanguage{english}
\textit{Ich versichere wahrheitsgemäß, die Arbeit selbstständig verfasst, alle benutzten Hilfsmittel vollständig und genau angegeben und alles kenntlich gemacht zu haben, was aus Arbeiten anderer unverändert oder mit Abänderungen entnommen wurde sowie die Satzung des KIT zur Sicherung guter wissenschaftlicher Praxis in der jeweils gültigen Fassung beachtet zu haben.}}
{\selectlanguage{ngerman}
\textit{Ich versichere wahrheitsgemäß, die Arbeit selbstständig verfasst, alle benutzten Hilfsmittel vollständig und genau angegeben und alles kenntlich gemacht zu haben, was aus Arbeiten anderer unverändert oder mit Abänderungen entnommen wurde sowie die Satzung des KIT zur Sicherung guter wissenschaftlicher Praxis in der jeweils gültigen Fassung beachtet zu haben.}}

\vspace{1cm}

%
%WICHTIG: Prüfen Sie unbedingt, ob der Text der Erklärung mit der jeweils für den Studiengang aktuell gültigen SPO übereinstimmt!
%
%TODO: Datum der Erklärung angeben
%            Hinweis: Das Datum der Erklärung ist üblicherweise das Abgabedatum
%
%TODO: Eigenen Vornamen und Nachnamen angeben
%
Karlsruhe, \today \hfill VORNAME NACHNAME

%\thispagestyle{empty}\cleardoublepage

%\mbox{}\thispagestyle{empty}


\end{document}